\scanpage{33}%
\mspage{18}\space
qu'il est divisible par $c_4^3-c_6^2$ dans $\mathbb{Z}/(2)[c_4,c_6,\Delta]$ (resp. $\mathbb{Z}/(3)[c_4,c_6,\Delta]$). Cela résulte de la \unsure{structure|structurc} de l'anneau des formes modulaires de caractéristique $2$ (resp. $3$).

\uline{Proposition} 7.2. \quad (I) \uline{L'anneau des formes modulaires de caractéristique} $2$ \uline{est} $\mathbb{Z}/(2)[\mbox{\begingroup\edef\encoutersep{\the\dimexpr.5\pgflinewidth\relax}\setbox0=\hbox{$a_1$}\setbox2=\hbox{\scriptsize\handwritten{le diviseur de $a_1$ est lisse sur $\mathbb{F}_2$}}\edef\encwda{\the\wd0}\edef\enchta{\the\ht0}\edef\encdpa{\the\dp0}\edef\encwdb{\the\wd2}\edef\enchtb{\the\ht2}\edef\encdpb{\the\dp2}\rlap{\smash{\tikz[baseline=0pt]{\begin{scope}[shift={(\the\dimexpr \encwda/2\relax,\the\dimexpr -\encdpa/2+\enchta/2\relax)}]\coordinate (enc-coeffa-north-east) at (\the\dimexpr \encoutersep+\encwda/2\relax,\the\dimexpr \encdpa/2+\enchta/2+\encoutersep\relax);\coordinate (enc-coeffa-south-west) at (\the\dimexpr -\encoutersep-\encwda/2\relax,\the\dimexpr -\encdpa/2-\enchta/2-\encoutersep\relax);\coordinate (enc-notea-west) at (\the\dimexpr -\encoutersep-\encwda/2+1.35em\relax,\the\dimexpr -\encdpa/2-\enchta/2-\encoutersep-1.05em\relax);\handwritten{\draw[thin] (enc-coeffa-south-west) to[out=-86,in=192] ([xshift=-.18em]enc-notea-west);}\pgfresetboundingbox\path[use as bounding box] ([yshift=-1.5em]enc-coeffa-south-west) rectangle (enc-coeffa-north-east);\end{scope}\pgfresetboundingbox\path[use as bounding box] (0pt,\the\dimexpr-\encdpa\relax) rectangle (\encwda,\enchta);}}}\rlap{\kern\the\dimexpr 1.35em\relax\raisebox{\the\dimexpr -\encdpa+\encdpb/2-\enchtb/2-\encoutersep-1.05em\relax}[0pt][0pt]{\copy2}}\copy0\endgroup},\Delta]$ , \uline{et} $c_4=a_1^4$ , $c_6=a_1^6$ .

\hspace*{3em}(II) \uline{L'anneau des formes modulaires de caractéristique} $3$ \uline{est} $\mathbb{Z}/(3)[\mbox{\begingroup\edef\encoutersep{\the\dimexpr.5\pgflinewidth\relax}\setbox0=\hbox{$b_2$}\setbox2=\hbox{\scriptsize\handwritten{le diviseur de $b_2$ est lisse sur $\mathbb{F}_2$}}\edef\encwda{\the\wd0}\edef\enchta{\the\ht0}\edef\encdpa{\the\dp0}\edef\encwdb{\the\wd2}\edef\enchtb{\the\ht2}\edef\encdpb{\the\dp2}\rlap{\smash{\tikz[baseline=0pt]{\begin{scope}[shift={(\the\dimexpr \encwda/2\relax,\the\dimexpr -\encdpa/2+\enchta/2\relax)}]\coordinate (enc-coeffb-north-east) at (\the\dimexpr \encoutersep+\encwda/2\relax,\the\dimexpr \encdpa/2+\enchta/2+\encoutersep\relax);\coordinate (enc-coeffb-south-west) at (\the\dimexpr -\encoutersep-\encwda/2\relax,\the\dimexpr -\encdpa/2-\enchta/2-\encoutersep\relax);\coordinate (enc-noteb-west) at (\the\dimexpr -\encoutersep-\encwda/2+1.35em\relax,\the\dimexpr -\encdpa/2-\enchta/2-\encoutersep-1.05em\relax);\handwritten{\draw[thin] (enc-coeffb-south-west) to[out=-86,in=192] ([xshift=-.18em]enc-noteb-west);}\pgfresetboundingbox\path[use as bounding box] ([yshift=-1.5em]enc-coeffb-south-west) rectangle (enc-coeffb-north-east);\end{scope}\pgfresetboundingbox\path[use as bounding box] (0pt,\the\dimexpr-\encdpa\relax) rectangle (\encwda,\enchta);}}}\rlap{\kern\the\dimexpr 1.35em\relax\raisebox{\the\dimexpr -\encdpa+\encdpb/2-\enchtb/2-\encoutersep-1.05em\relax}[0pt][0pt]{\copy2}}\copy0\endgroup},\Delta]$ , et $c_4=b_2^2$ , $c_6=-b_2^3$ .

Tout d'abord, il résulte de $(1.6)_1$ et $(1.7)_2$ que $a_1$ (resp. $b_2$) est une forme modulaire de caractéristique $2$ (resp. $3$) .

En caractéristique $2$ , cherchons tout d'abord les formes modulaires relatives aux seules courbes telles que $a_1$ soit inversible. Ces courbes se mettent sous forme (cf. (4.\mbox{\begingroup\edef\encoutersep{\the\dimexpr.5\pgflinewidth\relax}\setbox0=\hbox{1}\edef\encwda{\the\wd0}\edef\enchta{\the\ht0}\edef\encdpa{\the\dp0}\rlap{\smash{\tikz[baseline=0pt]{\begin{scope}[shift={(\the\dimexpr \encwda/2\relax,\the\dimexpr -\encdpa/2+\enchta/2\relax)}]\coordinate (enc-un-center) at (0pt,0pt);\handwritten{\draw[thin,gray,rotate around={-12:(enc-un-center)}] (enc-un-center) ellipse (0.36em and 1.12em);}\end{scope}\pgfresetboundingbox\path[use as bounding box] (0pt,\the\dimexpr-\encdpa\relax) rectangle (\encwda,\enchta);}}}\copy0\endgroup}))
\[
 y^2+a_1xy=x^3+a_2x^2+a_6 \qquad \pi=\frac{dx}{a_1x}
\]
et, compte tenu de (4.3), on voit que ces formes s'écrivent
\[
\begin{aligned}
 &a_1^{-n}\,P(a_1,a_6) \qquad \text{pour }P\text{ isobare, ou encore}\\[.65em]
 &a_1^{-m}\,Q(a_1,\Delta)\quad\handwritten{,}
\end{aligned}
\]
et cela garde un sens pour $a_1=0$ , $\Delta\ne0$ si et seulement si on avait en fait affaire à un polyn\mbox{\begingroup\edef\encoutersep{\the\dimexpr.5\pgflinewidth\relax}\setbox0=\hbox{\mbox{$\overset{\inserted{\handwritten{\wedge}}}{\mathrm{o}}$}}\edef\encwda{\the\wd0}\edef\enchta{\the\ht0}\edef\encdpa{\the\dp0}\rlap{\smash{\tikz[baseline=0pt]{\begin{scope}[shift={(\the\dimexpr \encwda/2\relax,\the\dimexpr -\encdpa/2+\enchta/2\relax)}]\coordinate (enc-oc-center) at (0pt,0pt);\handwritten{\draw[thin,gray,rotate around={-14:(enc-oc-center)}] (enc-oc-center) ellipse (0.56em and 1.1em);}\end{scope}\pgfresetboundingbox\path[use as bounding box] (0pt,\the\dimexpr-\encdpa\relax) rectangle (\encwda,\enchta);}}}\copy0\endgroup}me en $a_1$ et $\Delta$ .

En caractéristique $3$ , raisonnant de même d'abord sur les courbes telles que $c_4$ soit inversible, et remarquant que $d_6=0$ équivaut alors à $a_4=0$ dans (3.1), on voit que toute forme modu%
% archive sheet number: 32
